% TERMINALFIBER1 condensed, 2026-11-18 0445 KST
\documentclass[11pt]{article}
\usepackage{amsmath,amssymb,amsthm,hyperref}
\newtheorem{theorem}{Theorem}
\newtheorem{lemma}[theorem]{Lemma}
\title{TERMINALFIBER1: Terminal Output Multiplicity and Matched-Oracle PFCERT}
\author{P versus NP project, historic LSPV-5 certificate track}
\date{18 November 2026}
\begin{document}\maketitle
\noindent\textbf{Status:} The exact linear-output-fiber formula, scalable Reed--Muller moment obstruction, periodic affine endpoint fast-forward and matched-oracle two-output PFCERT impossibility are PROVED/EXACT in their stated models, pending external review. Finite tests are COMPUTATIONAL; unrestricted certificates and $P\ne NP$ are OPEN. This is a condensed public mirror; the complete local six-page report, executable tests and raw data are preserved separately.

\section{IDEAS PROPOSED BEFORE LITERATURE}
Before targeted primary-literature searching, the original four proposals and their falsifiers were frozen in \texttt{proposals\_before\_literature.md}, SHA-256 \texttt{07749d82fcf494fe99d24e52f9f1200092bb8bc77574217bd6452219308c5bb4}.

\noindent\textbf{CANDIDATE A.} Show that a degree-$s$ moment invariant on $2^d$ memory cells permits $2^k$ spurious ordinary readout labels for $k<2^{d-s}$. Predicted falsifier: a rank increment smaller than $k$.
\par\noindent\textbf{CANDIDATE B.} Prove the exact formula $|\{Lx+c:Hx=q\}|=2^{\operatorname{rk}[H;L]-\operatorname{rk}H}$ for consistent $q$ and use it as an explicit-matrix global-output test. Falsifier: exhaustive image-size mismatch.
\par\noindent\textbf{CANDIDATE C.} Fast-forward periodic affine transitions by repeated squaring to bind the actual reachable output after exponentially many steps. Falsifier: an endogenous branch violates the periodic schedule.
\par\noindent\textbf{CANDIDATE D.} Relativize Hirahara's $K^t$ immunity via one matched oracle and use a Baker--Gill--Solovay collapse to refute a universal two-label PFCERT. Falsifier: oracle mismatch, missing proof-size bound, or hidden descriptor selection.

\section{LITERATURE CHECK AND REVISIONS}
Hirahara, ECCC TR20-050 (2020), Theorem 3.9 and Section 9, proves $P$-immunity of $\{x:K^t(x)\ge |x|-1\}$ for superpolynomial $t$. Here $K^t$ is shortest description within time $t(|x|)$, not Levin's $\mathrm{Kt}(x)=\min_p(|p|+\log t_p)$. His prefix-plus-polynomial-suffix-search proof is explicitly relativized below with the SAME oracle. Baker--Gill--Solovay, SIAM J. Comput. 4(1975), constructs recursive oracle $A$ with $P^A=NP^A$. The minimum distance of the Boolean Reed--Muller code $RM(s,d)$ is $2^{d-s}$; see Saptharishi--Shpilka--Volk, arXiv:1503.09092. These are existing results; the precise PFCERT source-reachability application is new only as a restricted diagnostic. Source syntax, hash freshness or bitstream chronology alone does not show nonalgebrization.

\section{CONSTRUCTION OR PROOF ATTEMPT}
Fix binary $H\in\mathbb F_2^{r\times N}$, $L\in\mathbb F_2^{k\times N}$, $q\in\operatorname{im}H$ and $c\in\mathbb F_2^k$.
\begin{theorem}[Exact output fiber]
\[
\left|\{Lx+c:Hx=q\}\right|
=2^{\operatorname{rank}\binom HL-\operatorname{rank}H}.
\]
\end{theorem}
\begin{proof} Choose $x_0$ with $Hx_0=q$. The feasible fiber is $x_0+\ker H$, so its output set is an affine translate of $L(\ker H)$. Rank-nullity shows its dimension is $\operatorname{rank}\binom HL-\operatorname{rank}H$.
\end{proof}

Let $N=2^d$. Index cells by $a\in\{0,1\}^d$, and let $H_{d,s}$ have row $R\subseteq[d]$, $|R|\le s$, with entry $\prod_{i\in R}a_i$ in column $a$. Its row space is $RM(s,d)$, with $r=\sum_{j=0}^s\binom dj=O_s(d^s)$ rows.
\begin{theorem}[Scalable terminal ambiguity]
For any $k$ distinct ordinary cell-read probes, $1\le k<2^{d-s}$, and the associated coordinate map $L$,
\[
\operatorname{rank}\binom{H_{d,s}}L-\operatorname{rank}H_{d,s}=k,\qquad
\left|\{Lx:H_{d,s}x=0\}\right|=2^k.
\]
For fixed $s$ and all sufficiently large $d$, $k=d$ yields $2^d$ moment-compatible output strings although only $O_s(d^s)$ moments are stored.
\end{theorem}
\begin{proof}
Any nonzero linear combination of these $k$ coordinate rows has positive Hamming weight at most $k<2^{d-s}$. Every nonzero rowspace vector of $H_{d,s}$ has weight at least $2^{d-s}$ by Reed--Muller minimum distance. Thus the $k$ coordinates are linearly independent modulo this rowspace. Apply the fiber theorem.
\end{proof}
An actual ordinary deterministic program can initialize all $2^d$ tape cells to zero, perform optional exponential no-op delay, read $k=d$ cells and output $0^d$. Its true unique output is zero, whereas \emph{moment-only} terminal safety accepts all $2^d$ labels. The example refutes a moment-invariant-to-two-output inference, not a stronger source-semantic reachability certificate. It supplies no SAT lower bound.

\begin{theorem}[Positive restricted reachable endpoint]
For a fully source-authenticated, explicit period-$p$ affine update $z_{t+1}=A_{t\bmod p}z_t+b_{t\bmod p}$ on $m$ bits, and binary-encoded $T$, the exact reachable state after $T$ steps and every affine readout are computable in $O(pm^3+m^3\log(T+2)+\operatorname{poly}(pm^2+\log T))$ bit time and polynomial space. An exact source-output verifier therefore has globally ONE accepted output, despite $T$ potentially being exponential in the source size.
\end{theorem}
\begin{proof}
Embed $A_i,b_i$ in the $(m+1)$-square binary matrix $B_i=\left(\begin{smallmatrix}A_i&b_i\\0&1\end{smallmatrix}\right)$; set $M=B_{p-1}\cdots B_0$. For $T=qp+r$, apply $B_{r-1}\cdots B_0 M^q$ to $(z_0,1)$, with repeated squaring. All source parsing, matrix operations and output are explicitly charged. Endogenous schedule changes are outside this restricted grammar.
\end{proof}

Fix an oracle $O$, a universal oracle machine $U^O$ with polynomial-overhead universal simulation, and a superpolynomial time bound $t(n)$ that eventually dominates every polynomial. Set $R^{t,O}=\{x:K_{U}^{t,O}(x)\ge |x|-1\}$.
\begin{lemma}[Matched-oracle immunity]
$R^{t,O}$ is infinite and $P^O$-immune for every fixed legitimate oracle $O$.
\end{lemma}
\begin{proof} Counting programs shorter than $n-1$ leaves random strings of every length. If an infinite $P^O$ language $S$ were contained in it, fix a length $n$ having an $S$-string and hard-code the first $n-m$ bits of that string for $m=c\lceil\log_2 n\rceil$. Enumerate the $2^m$ suffixes, querying the SAME $P^O$ decider for $S$, and print its first accepted completion. This $O$-oracle program runs in polynomial time, eventually below $t(n)$, and has length at most $n-m+2\log n+O(1)<n-1$ for fixed sufficiently large $c$. Its output lies in $S$, a contradiction.
\end{proof}
\begin{theorem}[Collapsing-oracle PFCERT barrier]
For an oracle $A$ with $P^A=NP^A$, no polynomial-time $A$-oracle verifier $V^A(1^n,p,x,\pi)$ with $|\pi|\le\operatorname{poly}(n)$ can simultaneously satisfy (C) coverage of every length-$n$ output of $U^A(p)$ within $t(n)$ for $|p|\le n-2$, and (B2) at most TWO distinct accepted output labels $x$ over ALL proofs for every fixed $(n,p)$.
\end{theorem}
\begin{proof}
Let $E_n^A$ consist of all outputs accepted with some eligible $p,\pi$. As an existentially witnessed polynomial relation, $E^A\in NP^A$. There are $2^{n-1}-1$ eligible programs, so $|E_n^A|\le 2^n-2$. Thus $\overline{E^A}$ contains at least two strings of every length. Coverage places each excluded string in $R^{t,A}$. Under $NP^A=P^A$, $\overline{E^A}\in P^A$, contradicting matched-oracle immunity.
\end{proof}
This proves a necessary nonrelativization test only. It does NOT prove ordinary PFCERT exists; existence of the unrelativized universal schema would itself imply $P\ne NP$. The source grammar of the affine positive theorem is too restrictive to establish ordinary universal coverage.

\section{TESTS AND OUTCOME}
Executed \texttt{fiber\_rank\_audit.py}, revision \texttt{fiber-rank-audit.v1}, deterministic seeds 20261119--20261122, as four independent checkpointed runs. Batch 1: 20 rank-increment cases, 40 assertions, PASS. Batch 2: 11 code families and 67,840 minimum-distance assertions, PASS. Batch 3: 160 constructed label witnesses and brute-fiber checks, 162 assertions, PASS. Batch 4: 100 affine periodic fast-forward cases against independent iteration/cycle simulation, PASS. TOTAL: 291 cases, 68,142 assertions, zero failures. Exact code, per-case raw JSON, seed and SHA-256 manifest are in the local reproducibility bundle. Finite checks are COMPUTATIONAL, never asymptotic proofs.

\noindent\textbf{Four gates.} DISJ/hard-family: not a universal lower bound; affine remains an easy control. Relativization: local fiber and fast-forward lemmas relativize; universal PFCERT fails in a collapsing oracle. Simultaneous DTISP: polynomial time/space applies only to the fast-forward verifier, not to the original exponentially long run. Configuration counting: no lower bound from a polynomial moment summary; the $2^d$-cell tape can have $2^{2^d}$ configurations. No nonalgebrizing step, no Natural Proofs circuit property, no hash hardness.

\noindent\textbf{Cross-track ledger.} TMR-1 source-to-counting/hierarchy is conditional at the unreplayed kernel/NP bridge. TMR-2 exact-counting spectral algorithms are restricted. TMR-3 virtual checkpoint blank-start arbitrary provider remains open; recent research lives on draft PRs. TMR-4 costed shared-DAG-to-branching transfers have restricted valid source grammars and no universal size lower bound. TMR-5 full-chain red-team has restricted source-authenticated checks and unresolved universal proof transport. Current governance since November 15 assigns stable TMR roles; this historical LSPV-5 report did not edit active agent manuscripts, root main.tex, retired trackers or schedules.

\noindent\textbf{ONE next target.} Extend exact reachable-output certificates to a source-recognized class with endogenous nonlinear controller switches while proving an exact polynomial-cost quotient discovery and all terminal readouts; begin by finding the smallest switching counterexample.

\begin{thebibliography}{4}
\bibitem{hir} S. Hirahara, ECCC TR20-050 (2020), Theorem 3.9 and Section 9. \url{https://eccc.weizmann.ac.il/report/2020/050/download/}.
\bibitem{bgs} T. Baker, J. Gill and R. Solovay, SIAM J. Comput. 4(4), 431--442 (1975), DOI 10.1137/0204037.
\bibitem{rm} R. Saptharishi, A. Shpilka and B. L. Volk, arXiv:1503.09092 (2015).
\bibitem{baz} Y. Bazlov, Coding Theory Lecture Notes, Thm. 11.4, Manchester.
\end{thebibliography}
\end{document}
% FINAL_VERIFICATION_SENTINEL::LSPV5::TERMINALFIBER1::GITHUB_CONDENSED::2026-11-18_0445::P_VS_NP_OPEN
